\section{Results}
\label{sec:results}

\IfFileExists{numbers.tex}{}{%
\begin{center}
\fcolorbox{red!60!black}{red!5}{%
\begin{minipage}{0.94\textwidth}
\small\textbf{Status.} The pipeline has been validated end to end against
simulated injections (\cref{sec:validation}) and every data path has been
verified against the live services, but the production run over the full O3
and O4a samples has not yet been executed. Quantities awaiting that run are
shown as \result{like this}. They are placeholders, not measurements, and no
value in this section should be quoted until this box disappears --- it does
so automatically once \code{paper/numbers.tex} is written by stage~06 of the
pipeline. The structure below, the statistical treatment, and the caveats are
final.
\end{minipage}}
\end{center}
\vspace{0.6em}}

\subsection{The O3 sample}
\label{sec:results:sample}

\Cref{tab:sample} gives the size of the labelled sample after the cuts of
\cref{sec:data:gravityspy}.

\begin{table}[htbp]
\centering
\caption{Gravity Spy triggers surviving the selection cuts, by detector and
epoch, and the scattered-light fraction. Produced by
\code{results/tables/catalog\_summary.csv}.}
\label{tab:sample}
\small
\generated{tab_sample}
\end{table}

\begin{figure}[htbp]
\centering
\safefig[0.9\textwidth]{example_slow_1_L1}
\caption{A representative slow-scattering trigger with the extracted arch
envelope overlaid. The harmonic comb, the smooth rise and fall, and the
regular repetition are the features the analysis quantifies.}
\label{fig:exampleslow}
\end{figure}

\subsection{Morphology of the O3 population}
\label{sec:results:morphology}

The distribution of the arch repetition period is the central descriptive
result, because \cref{eq:fdrive} converts it directly into a statement about
the mechanical disturbance. \Cref{fig:archperiod} shows it with the implied
driving frequency on the upper axis.

\begin{figure}[htbp]
\centering
\safefig[0.9\textwidth]{arch_period_o3_vs_o4a}
\caption{Arch repetition period. The upper axis gives the implied driving
frequency $\fdrive = 1/2\Tarch$; the microseismic band
(\SIrange{0.1}{0.3}{\hertz}) and the anthropogenic band
(\SIrange{1}{6}{\hertz}) can be read off directly. The dotted line at
\SI{1.5}{\second} is the slow/fast split.}
\label{fig:archperiod}
\end{figure}

\Cref{tab:slowfast} decomposes each sample into the two populations. The
columns headed ``in band'' are the check that the decomposition means what it
claims: the slow population's implied driving frequencies should fall in the
microseismic band and the fast population's in the anthropogenic band. If they
do not, the split is a curve-fitting artefact rather than a physical
separation.

\begin{table}[htbp]
\centering
\caption{Slow/fast decomposition. ``In band'' is the fraction of each
population whose implied driving frequency falls in the corresponding
ground-motion band.}
\label{tab:slowfast}
\small
\generated{tab_slowfast}
\end{table}

\subsection{Auxiliary correlation in O3}
\label{sec:results:aux}

\subsubsection{Rate versus ground motion}

\Cref{tab:ratecorr} gives the Spearman correlation between the hourly
scattered-light rate and the hourly band-limited RMS of each ground-motion
channel, with six-hour moving-block bootstrap intervals.

\begin{table}[htbp]
\centering
\caption{Rank correlation of the hourly scattered-light rate with
ground-motion band-limited RMS, \Lone{} \Otzb. Channel names are abbreviated:
each is prefixed \code{L1:ISI-GND\_STS\_} and suffixed \code{.rms}. Intervals
are \SI{95}{\percent} moving-block bootstrap. Full table in
\code{results/tables/aux\_rate\_correlations.csv}.}
\label{tab:ratecorr}
\footnotesize
\generated{tab_ratecorr}
\end{table}

The strongest channel is \ResultStrongestChannel{}, at
$\rho = \ResultStrongestRho$.

Three supporting checks accompany it. The top-to-bottom quartile rate ratio
for that channel is \ResultQuartileRatio{}, which states the effect in
operational terms: how much more often scattering occurs in the noisiest
quarter of hours than in the quietest. The lag scan peaks at
\ResultLagPeak{}~hours, against the zero expected for an instantaneous
mechanical response. The partial correlation of rate with the anthropogenic
band, controlling for the microseism, tests whether the two bands carry
independent information or merely track one another.

\begin{figure}[htbp]
\centering
\begin{subfigure}{0.49\textwidth}
  \safefig[\textwidth]{aux_correlation_L1_O3b}
  \caption{Correlation summary.}
\end{subfigure}\hfill
\begin{subfigure}{0.49\textwidth}
  \safefig[\textwidth]{rate_vs_groundmotion_L1_O3b}
  \caption{Rate against ground motion.}
\end{subfigure}
\caption{Left: Spearman $\rho$ with moving-block bootstrap intervals for every
ground-motion channel. Right: hourly rate against the band-limited RMS of the
strongest channel, with binned medians and interquartile range.}
\label{fig:auxcorr}
\end{figure}

\subsubsection{The per-trigger predictor}

\Cref{tab:witness} gives the fitted witness calibration and the agreement
measured on the disjoint validation set. The calibration is quoted with the
range over which the training objective stays within \SI{1}{\percent} of its
maximum: a narrow range indicates that the witness genuinely constrains the
prediction, a wide one that it does not.

\begin{table}[htbp]
\centering
\caption{Scattering predictor $f_n = 2n|v|/\lambda$ against the witness
\code{L1:SUS-OM3\_M1\_DAMP\_L\_IN1\_DQ}. The calibration $\alpha$ is fitted on
a training set and frozen before the validation measurement; the
\SI{1}{\percent} range is the interval over which the training objective stays
within one per cent of its maximum.}
\label{tab:witness}
\small
\generated{tab_witness}
\end{table}

\begin{figure}[htbp]
\centering
\safefig[0.9\textwidth]{predictor_overlay_1_L1}
\caption{A real scattered-light trigger with the harmonic comb predicted from
the auxiliary witness overlaid (dashed). The prediction uses only auxiliary
data and the single fitted scale factor; the arches are from strain. Agreement
between the two is the evidence that \cref{eq:fn} describes the mechanism.}
\label{fig:predictor}
\end{figure}

\subsection{The O4a search}
\label{sec:results:o4a}

\Cref{tab:search} records the coverage of the livetime-limited search and the
number of candidates classified as scattered light. The identifier's
cross-validated performance on held-out O3 folds is
$\mathrm{AUC} = \ResultClassifierAuc$ with an efficiency of
\ResultClassifierEfficiency{} at a false-positive rate of $0.01$ against the
control classes --- with the two caveats of \cref{sec:methods:classifier}
attached.

\begin{table}[htbp]
\centering
\caption{Blind-search coverage and yield. Rates are per hour of \emph{analysed}
livetime; the coverage fraction is the analysed livetime as a fraction of the
epoch's total.}
\label{tab:search}
\small
\generated{tab_search}
\end{table}

\subsection{O3 versus O4a}
\label{sec:results:comparison}

\subsubsection{Rate}

With the identical search applied to both runs, the ratio of the \Ofa{} rate
to the O3 rate is \ResultRateRatioHOne{} at \Hone{} and \ResultRateRatioLOne{}
at \Lone, with Poisson intervals as given in \cref{tab:search}.

\subsubsection{Composition}

The headline structural question is whether the slow, microseism-driven
component was suppressed relative to the fast component. At \Hone{} the slow
fraction went from \ResultSlowFractionOThreeHOne{} in O3 to
\ResultSlowFractionOFouraHOne{} in \Ofa, a change of
\ResultSlowShiftHOne{}; at \Lone{} it went from
\ResultSlowFractionOThreeLOne{} to \ResultSlowFractionOFouraLOne{}, a change of
\ResultSlowShiftLOne{} (\SI{95}{\percent} Newcombe intervals).

\subsubsection{Morphology}

\Cref{tab:featurecomp} lists the features whose distributions differ most
between the runs, ranked by Cliff's $\delta$ rather than by $p$-value, for the
reason given in \cref{sec:methods:stats}. The largest effect is in
\ResultStrongestFeature{}, at $\delta = \ResultStrongestDelta$.

\begin{table}[htbp]
\centering
\caption{Feature-by-feature comparison, ranked by effect size. $p$-values are
Benjamini--Hochberg corrected across all twenty features. Full table in
\code{results/tables/feature\_comparison.csv}.}
\label{tab:featurecomp}
\small
\generated{tab_featurecomp}
\end{table}

\begin{figure}[htbp]
\centering
\safefig[\textwidth]{feature_distributions_o3_vs_o4a}
\caption{Morphology feature distributions for the O3 and O4a samples,
normalised to unit area. Both samples are produced by the same search and the
same classifier, so differences reflect the detector rather than the pipeline.}
\label{fig:featuredist}
\end{figure}

\subsection{Systematics}
\label{sec:results:systematics}

Four checks bound the result.

\paragraph{Label purity.} The Gravity Spy confidence cut is varied from
$0.80$ to $0.95$. The O3 morphology medians move by \result{shift}, which
bounds the effect of label contamination.

\paragraph{Trigger generator.} The two comparison modes of
\cref{sec:methods:search} --- like-for-like against the same search on O3, and
against the Gravity Spy sample --- are reported side by side. Agreement
indicates the result does not depend on how triggers were found; disagreement
is quantified and attributed.

\paragraph{Sampling.} The livetime-limited search covers only
\ResultCoverageOFoura{} of O4a, so it is blind to structure on timescales longer than the
sampling interval. The windows are drawn uniformly in livetime and stratified
across the run, which makes the sample representative of the run as a whole
but does not recover, for example, a two-week period of unusual behaviour.

\paragraph{Extraction bias.} The \SI{15}{\percent} high bias in maximum arch
frequency (\cref{sec:validation}) applies identically to both runs and so
cancels in the comparison. It is not corrected for, and no absolute statement
is made using that feature.
